\chapter{Laparoscopy}
\section*{What Is This Test?} Laparoscopy is a minimally invasive procedure that uses a camera through small abdominal incisions to examine or treat pelvic and abdominal disease.\section*{Alternative Names} Diagnostic laparoscopy; minimally invasive abdominal surgery; keyhole surgery.\section*{Principle} Direct visualization and instruments allow inspection, biopsy, fluid sampling, and selected treatment.\section*{Why This Test Is Ordered} It may evaluate unexplained abdominal or pelvic pain, endometriosis, infertility, masses, cancer spread, or peritoneal disease.\section*{Primary vs Secondary Test} It is an invasive secondary procedure after history, examination, imaging, and less-invasive tests.\section*{How This Test Is Performed} Under anesthesia, gas expands the abdomen and a camera and instruments are inserted through small ports.\section*{What Preparation Is Needed} Fasting, medicine instructions, and preoperative assessment are required; arrange transportation and recovery support.\section*{Understanding Results} Findings may include adhesions, endometriosis, inflammation, masses, fluid, or normal anatomy; biopsy results may take longer.\section*{Follow-up Tests} Pathology, cultures, imaging, or specialist treatment planning may follow.\section*{References} \begin{enumerate}\item MedlinePlus. Laparoscopy.\item Society of American Gastrointestinal and Endoscopic Surgeons. Laparoscopy guidance.\end{enumerate}\clearpage\section*{Understanding Results and Follow-up}
The laboratory report should identify the specimen, method, units, reference interval or decision limit, and whether the result is positive, negative, abnormal, or inconclusive. Reference intervals vary with age, sex, pregnancy, specimen type, assay, and laboratory; a result outside a range is not automatically a diagnosis, and a result within a range does not always exclude disease. Discuss unexpected results with the ordering clinician, who can decide whether repeat or confirmatory testing, treatment, or referral is appropriate. Do not change medicines or delay urgent care based on an isolated result.
\section*{Regulatory Status and Authorization Date}This chapter describes a test category, not a specific commercial product. FDA, CDSCO, EU IVDR, and MHRA approval or registration dates therefore cannot be assigned without the assay manufacturer, product name, instrument, and intended use. For laboratory-developed tests, an individual agency approval date may not exist. See the Preface for the interpretation of regulatory status.
\clearpage
